\section*{Hoofdstuk \ref{ch:g11:circuits-and-power} --- Elektrische circuits en vermogen}

\begin{solution}{exo:g11:circuits-and-power:1}
$q = It = 2.0 \times 90 \times 60 = \qty{1.1e4}{C}$. Aantal elementaire
ladingen: $q/e = 1.08\times10^{4}/1.6\times10^{-19} \approx \num{6.8e22}$.
\end{solution}

\begin{solution}{exo:g11:circuits-and-power:2}
$I = U/R = 12/220 = \qty{0.055}{A} = \qty{55}{mA}$;
$R = U/I = 230/0.50 = \qty{460}{\ohm}$.
\end{solution}

\begin{solution}{exo:g11:circuits-and-power:3}
Serie: $120 + 60 = \qty{180}{\ohm}$. Parallel: $\frac{120 \times 60}{180}
= \qty{40}{\ohm}$ --- kleiner dan beide: zelfde
\omterm{def:g11:circuits-and-power:voltage}{spanning}, maar de stroom krijgt
een tweede pad, dus meer totale stroom en minder
\omterm{def:g11:circuits-and-power:resistance}{weerstand}.
\end{solution}

\begin{solution}{exo:g11:circuits-and-power:4}
$P = UI = 230 \times 8.7 = \qty{2.0}{kW}$. In $t = \qty{240}{s}$:
$E = Pt = 2000 \times 240 = \qty{4.8e5}{J} = 4.8\times10^{5}/3.6
\times10^{6} \approx \qty{0.13}{kW.h}$.
\end{solution}

\begin{solution}{exo:g11:circuits-and-power:5}
$E = 60 \times 5.0 = \qty{300}{W.h} = \qty{0.30}{kW.h} = 0.30 \times
3.6\times10^{6} = \qty{1.1e6}{J}$. Kostprijs: $0.30 \times 0.25 = 0.075$
euro per dag.
\end{solution}

\begin{solution}{exo:g11:circuits-and-power:6}
\emph{1.} $R = U^2/P = 230^2/1200 = \qty{44}{\ohm}$;
$I = P/U = 1200/230 = \qty{5.2}{A}$.

\emph{2.} $P' = U'^2/R = 115^2/44.1 = \qty{300}{W}$: een kwart, niet
de helft, omdat halveren van $U$ ook $I$ halveert, en $P = UI \propto
U^2$.
\end{solution}

\begin{solution}{exo:g11:circuits-and-power:7}
$R_{\text{s}} = \qty{60}{\ohm}$, dus $I = 12/60 = \qty{0.20}{A}$.
$U_1 = R_1 I = \qty{8.0}{V}$, $U_2 = \qty{4.0}{V}$ (som \qty{12}{V});
$P_1 = U_1 I = \qty{1.6}{W}$, $P_2 = \qty{0.8}{W}$: totaal \qty{2.4}{W}
$= UI = 12 \times 0.20$, zoals het moet.
\end{solution}

\begin{solution}{exo:g11:circuits-and-power:8}
$I_1 = 12/60 = \qty{0.20}{A}$, $I_2 = 12/30 = \qty{0.40}{A}$, totaal
$I = \qty{0.60}{A}$. Equivalent: $R = U/I = 12/0.60 = \qty{20}{\ohm}$,
en inderdaad $\frac{60 \times 30}{60 + 30} = \qty{20}{\ohm}$.
\end{solution}

\begin{solution}{exo:g11:circuits-and-power:9}
$\mathcal{E} - 0.50\,r = 5.5$ en $\mathcal{E} - 2.0\,r = 4.0$;
aftrekken, $1.5\,r = 1.5$, dus $r = \qty{1.0}{\ohm}$ en $\mathcal{E} =
\qty{6.0}{V}$. Kortsluiting: $\mathcal{E}/r = \qty{6.0}{A}$.
\end{solution}

\begin{solution}{exo:g11:circuits-and-power:10}
Waterkoker $2200/230 = \qty{9.6}{A}$, toaster $1000/230 = \qty{4.3}{A}$:
totaal \qty{13.9}{A} $< \qty{16}{A}$, de
\omterm{def:g11:circuits-and-power:joule}{zekering} houdt. Verwarming
$1500/230 = \qty{6.5}{A}$: totaal \qty{20.4}{A} $> \qty{16}{A}$, ze
doorslaat (equivalent, $\qty{4.7}{kW} > 230 \times 16 \approx
\qty{3.7}{kW}$).
\end{solution}

\begin{solution}{exo:g11:circuits-and-power:11}
Bij \qty{200}{V}: $I = 10^4/200 = \qty{50}{A}$, verlies $RI^2 = 2.0
\times 50^2 = \qty{5.0}{kW}$ --- de helft van de levering. Bij
\qty{4.0}{kV}: $I = \qty{2.5}{A}$, verlies $2.0 \times 2.5^2 =
\qty{12.5}{W}$. Verhouding $400 = 20^2$: twintig keer de
\omterm{def:g11:circuits-and-power:voltage}{spanning} betekent twintig keer
minder stroom, en het verlies gaat als $I^2$.
\end{solution}

\begin{solution}{exo:g11:circuits-and-power:12}
Vier lay-outs: alle in serie, $3 \times 60 = \qty{180}{\ohm}$; alle
parallel, $60/3 = \qty{20}{\ohm}$; één in serie met twee parallel,
$60 + 30 = \qty{90}{\ohm}$; één parallel met twee in serie,
$\frac{60 \times 120}{180} = \qty{40}{\ohm}$.
\end{solution}

\begin{solution}{exo:g11:circuits-and-power:13}
\emph{1.} $I = \mathcal{E}/(R + r) = 4.5/7.5 = \qty{0.60}{A}$;
$U = RI = \qty{3.6}{V}$; nuttig $P = UI = \qty{2.2}{W}$; verloren
$rI^2 = 1.5 \times 0.60^2 = \qty{0.54}{W}$; $\eta = U/\mathcal{E} =
3.6/4.5 = 0.80$.

\emph{2.} Volgens de hint, $P = \mathcal{E}^2 \big/ \big[(R - r)^2/R +
4r\big]$; de haak is minimaal ($= 4r$) precies wanneer $R = r$, dus
$P_{\max} = \mathcal{E}^2/(4r) = 4.5^2/6.0 = \qty{3.4}{W}$.
\end{solution}

\begin{solution}{exo:g11:circuits-and-power:14}
Warmte: $Q = 1.0 \times 4180 \times (100 - 20) = \qty{3.3e5}{J}$.
Elektrisch: $E = Q/0.90 = \qty{3.7e5}{J}$. Tijd: $t = E/P =
3.7\times10^{5}/2200 \approx \qty{170}{s}$, onder drie minuten.
Kostprijs: $3.7\times10^{5}/3.6\times10^{6} \approx \qty{0.10}{kW.h}$,
ongeveer $2.6$ cent.
\end{solution}

\begin{solution}{exo:g11:circuits-and-power:15}
\emph{1.} $q = 60 \times 3600 = \qty{2.2e5}{C}$; $E \approx \mathcal{E}q
= 12.7 \times 2.16\times10^{5} = \qty{2.7e6}{J} \approx \qty{0.76}{kW.h}$.

\emph{2.} $U = 12.7 - 0.020 \times 150 = \qty{9.7}{V}$; geleverd
$UI = \qty{1.5}{kW}$; verloren $rI^2 = 0.020 \times 150^2 = \qty{450}{W}$.

\emph{3.} De koplampen zitten over de klemmen: tijdens het starten daalt
de klem-\omterm{def:g11:circuits-and-power:voltage}{spanning} van $12.7$ naar
$\qty{9.7}{V}$, dus ze ontvangen minder
\omterm{def:g10:energy-conservation:power}{vermogen} en dimmen.
\end{solution}

\begin{solution}{pb:g11:circuits-and-power:1}
\textbf{1.} Elke seconde, $E_p = mgh = 5.0\times10^{4} \times 9.81 \times
200 = \qty{9.8e7}{J}$: beschikbaar
\omterm{def:g10:energy-conservation:power}{vermogen} \qty{98}{MW}.

\textbf{2.} $0.90 \times 98 = \qty{88}{MW}$.

\textbf{3.} $I = P/U = 8.8\times10^{7}/2.0\times10^{4} = \qty{4.4e3}{A}$.

\textbf{4.} $R_\ell I^2 = 32 \times (4.4\times10^{3})^2 \approx
\qty{6.2e8}{W} = \qty{620}{MW}$ --- zeven keer de output van de centrale:
transmissie bij \qty{20}{kV} is onmogelijk.

\textbf{5.} Met $P = UI$ vast, verkleint alleen $U$ verhogen $I$: een
transformator.

\textbf{6.} $I = 8.8\times10^{7}/4.0\times10^{5} = \qty{220}{A}$.

\textbf{7.} $R_\ell I^2 = 32 \times 220^2 \approx \qty{1.6}{MW}$.

\textbf{8.} $1.6/88 \approx 0.018$: ongeveer $1.8\%$.

\textbf{9.} Fractie $= R_\ell I^2/P = R_\ell (P/U)^2/P = R_\ell P/U^2
\propto 1/U^2$: van \qty{20}{kV} naar \qty{400}{kV} daalt ze met
$20^2 = 400$, en inderdaad $\qty{620}{MW}/400 \approx \qty{1.6}{MW}$.

\textbf{10.} $R_\ell I = 32 \times 220 \approx \qty{7.1}{kV}$: $1.8\%$
van \qty{400}{kV}, dezelfde fractie als vraag~8, want
$R_\ell I/U = R_\ell I^2/UI$.

\textbf{11.} $\qty{1}{kW.h} = \qty{3.6e6}{J}$: de meter vermenigvuldigt
het getrokken \omterm{def:g10:energy-conservation:power}{vermogen} met de
tijd dat het getrokken wordt, en telt op.

\textbf{12.} $I = 900/230 = \qty{3.9}{A}$;
$R = U^2/P = 230^2/900 \approx \qty{59}{\ohm}$.

\textbf{13.} $E = 900 \times 180 = \qty{1.6e5}{J} = \qty{0.045}{kW.h}$:
ongeveer $1.1$ cent.

\textbf{14.} De fysica van $RI^2$ is identiek; alleen het doel
verschilt --- op de lijn is de warmte ongewenst, in de toaster \emph{is}
de warmte het product.

\textbf{15.} $3.9 + 9.6 = \qty{13.5}{A} < \qty{16}{A}$: houdt. Verwarming
$2000/230 = \qty{8.7}{A}$: totaal $\qty{22.2}{A} > \qty{16}{A}$, de
\omterm{def:g11:circuits-and-power:joule}{zekering} doorslaat.

\textbf{16.} $0.05 \times 13.5^2 = \qty{9.1}{W}$: $9.1/3100 \approx
0.3\%$ van het geleverde
\omterm{def:g10:energy-conservation:power}{vermogen} --- verwaarloosbaar over
korte draden.

\textbf{17.} $\eta = 0.90 \times 0.99 \times 0.982 \times 0.99 \times
0.997 \approx 0.86$.

\textbf{18.} Eén getoast
\omterm{def:g10:energy-conservation:kwh}{kilowattuur} heeft
$3.6\times10^{6}/0.86 \approx \qty{4.2e6}{J}$ stroomopwaarts nodig:
$m = E/(gh) = 4.2\times10^{6}/(9.81 \times 200) \approx \qty{2.1e3}{kg}$,
ongeveer \qty{2.1}{m^3} water.

\textbf{19.} $R_\ell P/U^2 = 32 \times 8.8\times10^{7}/(2.0\times
10^{4})^2 \approx 7$. Een ``fractie'' boven $1$ betekent dat de lijn
meer zou dissiperen dan ze voert: de levering kan bij die
\omterm{def:g11:circuits-and-power:voltage}{spanning} helemaal niet
gebeuren.

\textbf{20.} Eén getoast
\omterm{def:g10:energy-conservation:kwh}{kilowattuur} kost de berg ongeveer
twee ton --- \qty{2.1}{m^3} --- water dat \qty{200}{m} valt, en zo'n
$14\%$ van de \omterm{def:g10:energy-conservation:energy}{energie} sterft
onderweg, vooral in turbines en transformatoren (onder $2\%$ in
\qty{1000}{km} lijn). De enige truc die de reis mogelijk maakt is de
transformator: omhoogtransformeren naar \qty{400}{kV} deelt de stroom
door $20$ en het joule-verlies door $400$.
\end{solution}
